一半径为 R、质量为 M 的均匀密度球形天体静止，万有引力常数 G。沿距离球心为 d∈[0,R) 的直线挖出质量可忽略的狭窄真空隧道，隧道两端在球面。质量 m 的微小滑块从隧道一端静止释放，在无摩擦轨道内沿隧道运动，天体不自转、忽略滑块反作用。
(1) 由壳层定理求球内引力场和势能；(2) 设 s 为从弦中点起的带符号坐标，求精确运动方程与到达另一端所需时间；(3) 求最大速度及沿轨道各点的速度函数；(4) 检验 d=0 与 d→R⁻ 极限并说明为何所有不同斜弦的穿越时间相同。
